\setlength{\marginparwidth}{4.5cm}

\usepackage[utf8]{inputenc}
%\usepackage[T1]{fontenc}
\usepackage{etex}

\pdfoutput=1
%\usepackage{subcaption}
\usepackage[english]{babel}
%\usepackage{wrapfig}
\let\proof\relax\let\endproof\relax
\usepackage{amsthm}
\usepackage{hyperref}
\usepackage{color}
\usepackage{mathtools}
% \usepackage{amsmath} loaded by mathtools
\usepackage{amssymb}
\renewcommand{\qedsymbol}{$\blacksquare$}
\usepackage{amsfonts}
%\usepackage{quantikz}
\usepackage{authblk}
\renewcommand\Authfont{\normalsize}

\usepackage{physics}
\usepackage{float}
\usepackage{enumitem}
\setlist{noitemsep} %no separation between list items
%\usepackage{algcompatible}
\usepackage{algpseudocodex}
\algrenewcommand\algorithmicprocedure{$\triangleright$}

%\usepackage{subfigure}
\usepackage{graphicx}
\usepackage{subcaption}
%\usepackage[dvipsnames]{xcolor}
\usepackage{bm} 
%\usepackage{qcircuit}
\usepackage{adjustbox}
\usepackage{multirow}
%% \ifcomment
%% \else
%\usepackage[table,xcdraw]{xcolor}
%% \fi % xcolor is already loaded by the tufte-handout class
%\usepackage{tikz, graphics} % graphics, if the figure contains embedded images
%\usetikzlibrary{arrows.meta, % if the figure contains arrow-tips
%                bending,     % arrow tips on arcs are "bent," i.e., deformed a bit
%                patterns,     % if the figure contains pattern fills
%                shapes,positioning
%               }

%\usetikzlibrary{arrows.meta,      % if the figure contains arrow-tips
%	bending,          % arrow tips on arcs are "bent," i.e., deformed a bit
%	patterns,         % if the figure contains pattern fills
%	shapes,
%	positioning}
% tikzstyles

%\tikzstyle{node}=[draw=black, minimum size = 10mm, shape = circle]
%\tikzstyle{nodep}=[draw=black, minimum size = 5mm, shape = circle]
%\tikzstyle{vertex}=[draw=black, fill = black, shape = circle, scale = .5]
%\tikzstyle{snode}=[draw=black, shape = circle]
%\tikzstyle{onode}=[draw=black, minimum size = 5mm, shape = rectangle]
%\tikzstyle{none}=[]
%\tikzstyle{dashed edge}=[-, dashed, dash pattern=on 2pt off 1.5pt, thick, draw={rgb,255: red,68; green,136; blue,255}]

% https://tex.stackexchange.com/questions/75836/creating-a-seamless-xor-symbol-as-node
%\tikzset{XOR/.style={draw,circle,append after command={
%			[shorten >=\pgflinewidth, shorten <=\pgflinewidth,]
%			(\tikzlastnode.north) edge (\tikzlastnode.south)
%			(\tikzlastnode.east) edge (\tikzlastnode.west)
%		}
%	}
%}

%\tikzset{bullet/.style={circle, fill,minimum size=4pt,inner sep=0pt, outer sep=0pt}}


\usepackage [
  n,
  advantage,
  operators,
  sets,
  adversary,
  landau,
  probability, 
  notions,
  logic,
  ff,
  mm,
  primitives,
  events,
  complexity,
  asymptotics,
  keys
  ] {cryptocode}
 

\usepackage[strings]{underscore}
\usepackage{bm}
\usepackage{booktabs} % nicer tables

\renewcommand{\cc}[1]{\complclass{#1}}
\newcommand{\Abort}{\mathsf{Abort}}
\newcommand{\Acc}{\mathsf{Accept}}
\newcommand{\Enc}{\mathsf{Enc}}
\newcommand{\Dec}{\mathsf{Dec}}
\newcommand{\Redo}{\mathsf{Redo}}
\newcommand{\BQP}{\mathsf{BQP}}
\newcommand{\BPP}{\mathsf{BPP}}

\newcommand{\X}{\mathsf{X}}
\newcommand{\Y}{\mathsf{Y}}
\newcommand{\Z}{\mathsf{Z}}
\newcommand{\Id}{\mathsf{I}}
\newcommand{\Ha}{\mathsf{H}}
\newcommand{\CZ}{\mathsf{CZ}}
\newcommand{\D}{\mathsf{D}}
\newcommand{\cZ}{\mathsf{CZ}}
\newcommand{\cX}{\mathsf{CX}}
\newcommand{\Sp}{\mathsf{S}}
\newcommand{\CNOT}{\mathsf{CNOT}}
\newcommand{\QOTP}{\mathsf{QOTP}}
\newcommand{\T}{\mathsf{T}}
\newcommand{\W}{\mathsf{W}}


\newcommand{\cptp}[1]{\mathsf{#1}} % CPTP notation
\newcommand{\pd}[1]{\mathcal{#1}}  % Probability distribution notation
\newcommand{\sch}[1]{\bm{#1}} % scheme


%\floatstyle{ruled}

\usepackage[capitalise]{cleveref}
 % \let\oldcref\cref
 % \renewcommand\cref[1]{\oldcref{#1}\textsuperscript{#1}}


\usepackage{thmtools}
\declaretheorem[name=Protocol, refname={Protocol, Protocols}]{protocol}
\declaretheorem[name=Resource, refname={Resource, Resources}]{resource}
\declaretheorem[name=Routine, refname={Routine, Routines}]{routine}

% \newfloat{simulator}{htb!}{idf}
% \floatname{simulator}{Simulator}

% \newfloat{routine}{htb!}{idf}
% \floatname{routine}{Routine}

% \ifcomment
% %% \setcounter{secnumdepth}{2}
% %% \newtheorem{theorem}{Theorem}
% %% \newtheorem{lemma}[theorem]{Lemma}
% %% \newtheorem{definition}[theorem]{Definition}
% %% \newtheorem{corollary}[theorem]{Corollary}
% %% \newtheorem{remark}[theorem]{Remark}
% %% \newcommand{\comment}[2]{\sidenote{{\textbf{#1}: #2}}}
% %% \newcommand{\mn}[2]{\textsuperscript{\textcolor{red}{\textsf{\textbf{#1}}}}\marginnote{\tiny \textcolor{red}{\textsf{{\textbf{#1:}} #2}}}}


%\usepackage{marginnote}
% \newcommand{\comment}[2]{\textsuperscript{\textcolor{red}{\textsf{\textbf{#1}}}}\marginnote{\tiny {\textsf{{\textbf{#1:}} #2}}}}
% \newcommand{\pn}[2]{\textcolor{red}{\textsf{\textbf{#1:} #2}}}
% \else
% \newcommand{\mn}[2]{}
% \newcommand{\pn}[2]{}
% \newcommand{\comment}[2]{}
% \fi


\usepackage[a4paper,lmargin=1.25in, rmargin=1.25in]{geometry}

%\usepackage[disable]{todonotes} % use to remove notes
\usepackage{todonotes}
\todostyle{holl}{color=green!30, size=\footnotesize}
\todostyle{as}{color=blue!30, size=\footnotesize}

\setcounter{secnumdepth}{3}

% \usepackage{markdown}
% \usepackage[hashEnumerators,texMathDollars]{markdown}

\theoremstyle{plain}
\newtheorem{theorem}{Theorem}[section]
\newtheorem{lemma}[theorem]{Lemma}
\newtheorem{corollary}[theorem]{Corollary}

\theoremstyle{definition}
\newtheorem{definition}[theorem]{Definition}
% \newtheorem{resourceenv}{Resource}
% \newtheorem{protocolenv}{Protocol}
% \newtheorem{simulatorenv}{Simulator}
